\subsection{Emparejamiento Bipartito Máximo, Kuhn (parejas de baile)}
\label{alg:4-1}

\graybold{Preguntas guía:}

\begin{itemize}
\item ¿Tengo dos conjuntos disjuntos (personas / tareas) y aristas que solo conectan un conjunto con el otro?
\item ¿Me piden el número máximo de "parejas" o asignaciones simultáneas sin repetir elementos de ningún lado?
\item ¿El tamaño de ambos lados permite \bigO{V \cdot E} (hasta unos pocos miles)?
\end{itemize}

\graybold{Utilidad:} Encontrar el máximo número de aristas seleccionables en un grafo bipartito sin que dos aristas compartan vértice (asignación óptima uno a uno). Ejemplos: asignar trabajadores a tareas, parejas de baile, cobertura de tareas.

\graybold{Complejidad:} \bigO{V \cdot E} con la versión simple de caminos aumentantes, donde V es el lado izquierdo y E el número de aristas.

\graybold{Fundamento teórico:} Se buscan repetidamente "caminos aumentantes": para cada vértice izquierdo se intenta, vía DFS, encontrar una pareja libre o "desalojar" a otro vértice ya emparejado para liberar espacio y aumentar el emparejamiento en 1. Es un caso particular de Ford-Fulkerson: caminos aumentantes en la red de flujo con capacidades unitarias. Por el teorema de Kőnig, el tamaño del emparejamiento es también el de la cobertura mínima de vértices.~\cite{kuhn1955,hopcroft1973}

\graybold{Ejercicio aplicado}

Hay n chicos y m chicas y k parejas de baile posibles. Hallar el número máximo de parejas simultáneas y mostrar una asignación válida.

\graybold{I/O:} n, m ≤ 500 y k ≤ 1000 → entrada pequeña, se usa readline por línea (patrón 1), suficiente y más simple de leer.

\graybold{Código solución}

\begin{lstlisting}[language=Python]
import sys
input = sys.stdin.buffer.readline
 
def try_kuhn(u, adj, matched, visited):
    # intenta emparejar el vertice u del lado izquierdo
    for v in adj[u]:
        if not visited[v]:
            visited[v] = True
            # v esta libre, o su pareja actual puede reubicarse...
            if matched[v] == -1 or try_kuhn(matched[v], adj, matched, visited):
                matched[v] = u
                return True
    return False
 
def solve():
    n, m, k = map(int, input().split())
    adj = [[] for _ in range(n)]
    for _ in range(k):
        a, b = map(int, input().split())
        adj[a - 1].append(b - 1)
 
    matched = [-1] * m                  # pareja actual de cada vertice derecho
    total = 0
    for u in range(n):
        visited = [False] * m           # se reinicia en cada intento
        if try_kuhn(u, adj, matched, visited):
            total += 1
 
    out = [str(total)]
    for v in range(m):
        if matched[v] != -1:
            out.append(f"{matched[v] + 1} {v + 1}")
    sys.stdout.write("\n".join(out))
 
solve()
\end{lstlisting}
